\section{Allgemeine Boxen}

\subsection{Die Umgebung \texttt{FlexibleBox}}

Die Umgebung \verb|FlexibleBox| ist die grundlegende Box-Umgebung der Klasse.

Sie ermöglicht es, beliebige Inhalte mit einem Titel und frei wählbaren Farben darzustellen.

Viele andere Befehle der Klasse, beispielsweise \verb|\sammelprobe| und \verb|Wertekasten|, verwenden intern diese Umgebung.

\subsubsection*{Syntax}

\begin{Verbatim}
\begin{FlexibleBox}
    {Titel}
    [Hintergrundfarbe]
    [Titelfarbe]
    [Textfarbe]

Inhalt

\end{FlexibleBox}
\end{Verbatim}

\subsubsection*{Parameter}

\begin{center}
\begin{tabular}{|l|l|l|}
\hline
Parameter & Bedeutung & Standardwert\\
\hline
Titel & Überschrift der Box & --\\
\hline
Hintergrundfarbe & Hintergrund des Inhalts & PergamentHell\\
\hline
Titelfarbe & Hintergrund des Titels & PergamentDunkel\\
\hline
Textfarbe & Farbe der Titelbeschriftung & white\\
\hline
\end{tabular}
\end{center}

Die Farben können alle von \LaTeX\ bekannten Farben sein, beispielsweise
\verb|yellow|, \verb|red| oder eigene mit \verb|\definecolor| definierte Farben.

\subsubsection*{Beispiel 1: Standardfarben}

\paragraph{Quellcode}

\begin{Verbatim}
\begin{FlexibleBox}{Hinweis}

Der alte Turm steht seit Jahrhunderten
verlassen im Wald.

\end{FlexibleBox}
\end{Verbatim}

\paragraph{Ergebnis}

\begin{FlexibleBox}{Hinweis}

Der alte Turm steht seit Jahrhunderten
verlassen im Wald.

\end{FlexibleBox}

\subsubsection*{Beispiel 2: Eigene Farben}

\paragraph{Quellcode}

\begin{Verbatim}
\begin{FlexibleBox}
    {Warnung}
    [yellow]
    [red]
    [white]

Die Tür ist magisch versiegelt.

\end{FlexibleBox}
\end{Verbatim}

\paragraph{Ergebnis}

\begin{FlexibleBox}
    {Warnung}
    [yellow]
    [red]
    [white]

Die Tür ist magisch versiegelt.

\end{FlexibleBox}
